\documentclass[a4paper,12pt]{article}

\usepackage[T2A]{fontenc}
\usepackage[utf8]{inputenc}
\usepackage[russian]{babel}

\usepackage{amsmath,amssymb,amsthm,mathtools}

% Sans-serif, совместимый с кириллицей T2A
\renewcommand{\familydefault}{\sfdefault}

\usepackage{icomma}
\usepackage{geometry}
\usepackage[
  protrusion=true,
  expansion=false,
  nopatch=footnote
]{microtype}

\usepackage{tikz}
\usetikzlibrary{calc}
\usetikzlibrary{arrows.meta,positioning,shapes.geometric}

\usepackage{tkz-euclide}

\usepackage{pgfplots}
\pgfplotsset{compat=1.18}

\usepackage{tabularx,booktabs,longtable,array}
\usepackage{enumitem}
\usepackage{hyperref}
\usepackage{parskip}
\usepackage{fancyhdr}
\usepackage{setspace}
\usepackage{xcolor}

\geometry{
  left=20mm,
  right=20mm,
  top=20mm,
  bottom=22mm
}

\onehalfspacing

\setlength{\parindent}{0pt}
\setlength{\parskip}{0.55em}
\setlength{\LTpre}{0pt}
\setlength{\LTpost}{0pt}

\setlist[itemize]{
  leftmargin=1.4em,
  itemsep=0.25em,
  topsep=0.25em
}

\setlist[enumerate]{
  leftmargin=1.7em,
  itemsep=0.45em,
  topsep=0.35em
}

\renewcommand{\arraystretch}{1.35}

% ============================================================
% Палитра
% ============================================================

\definecolor{accent}{HTML}{008080}
\definecolor{darkaccent}{HTML}{004D4D}
\definecolor{textgray}{HTML}{333333}
\definecolor{softgray}{HTML}{6B7280}
\definecolor{lightline}{HTML}{D8E5E5}

\color{textgray}

\hypersetup{
  colorlinks=true,
  linkcolor=darkaccent,
  urlcolor=darkaccent
}

% ============================================================
% Колонтитулы
% ============================================================

\pagestyle{fancy}
\fancyhf{}

\fancyhead[L]{\small МЦКО · геометрия · 7 класс}
\fancyhead[R]{\small Екатерина}
\fancyfoot[C]{\small \thepage}

\renewcommand{\headrulewidth}{0.3pt}
\renewcommand{\footrulewidth}{0pt}

\setlength{\headheight}{14.5pt}

% ============================================================
% tkz-euclide
% ============================================================

\tkzSetUpLine[line width=1pt,color=accent]
\tkzSetUpPoint[color=accent,fill=accent,size=3]
\tkzSetUpLabel[color=black,font=\small]

% ============================================================
% Стили блок-схем
% ============================================================

\tikzset{
  flowstart/.style={
    ellipse,
    draw=darkaccent,
    line width=0.9pt,
    minimum width=34mm,
    minimum height=10mm,
    align=center,
    inner sep=4pt
  },
  flowact/.style={
    rectangle,
    rounded corners=2mm,
    draw=darkaccent,
    line width=0.9pt,
    text width=77mm,
    minimum height=11mm,
    align=center,
    inner sep=5pt
  },
  flowcond/.style={
    diamond,
    aspect=2.2,
    draw=darkaccent,
    line width=0.9pt,
    text width=45mm,
    align=center,
    inner sep=2pt
  },
  flowarr/.style={
    -{Latex[length=2.4mm]},
    line width=0.9pt,
    draw=darkaccent
  },
  flowlabel/.style={
    font=\small,
    fill=white,
    inner sep=1.5pt
  }
}

% ============================================================
% Служебные команды
% ============================================================

\newcommand{\sectionline}{
  \par
  \vspace{0.2em}
  \textcolor{lightline}{\rule{\linewidth}{0.8pt}}
  \par
  \vspace{0.35em}
}

\newcommand{\key}[1]{
  \textcolor{darkaccent}{\textbf{#1}}
}

\newcommand{\checkitem}{
  \textcolor{accent}{$\square$}\,
}

\begin{document}

% ============================================================
% ТИТУЛЬНЫЙ ЛИСТ
% ============================================================

\thispagestyle{empty}

\begin{center}

\vspace*{23mm}

{\Large
\textcolor{darkaccent}{\textbf{Чек-лист}}
}
\par

\vspace{7mm}

{\huge
\textbf{Углы, треугольники и признаки равенства}
}
\par

\vspace{4mm}

{\Large
Подготовка к МЦКО · 7 класс
}
\par

\vspace{16mm}

{\large Екатерина}
\par

\vspace{4mm}

{\large \today}
\par

\vfill

{\large
Лёвин Артём Александрович
}
\par

\vspace{2mm}

{\normalsize
эксклюзивно для Екатерины
}
\par

\vspace{13mm}

\begin{tikzpicture}[x=1cm,y=1cm]
\draw[
  accent,
  line width=1.1pt
]
(-6,0)
.. controls (-3.6,0.9) and (-1.7,-0.8) ..
(0,0)
.. controls (1.8,0.8) and (3.7,-0.7) ..
(6,0);
\end{tikzpicture}

\vspace{8mm}

\end{center}

\clearpage

% ============================================================
% КАРТА ТЕМЫ
% ============================================================

\section*{Карта темы}
\sectionline

После повторения Вы должны уверенно:

\begin{itemize}

\item
различать \key{вертикальные} и \key{смежные} углы;

\item
применять свойства углов при
\key{параллельных прямых и секущей};

\item
использовать сумму углов треугольника и свойства
\key{равнобедренного треугольника};

\item
узнавать три признака равенства треугольников:
\key{СУС, УСУ, ССС};

\item
правильно устанавливать
\key{соответствие вершин}
равных треугольников;

\item
доказывать равенство отрезков и углов
через равенство подходящих треугольников;

\item
замечать недостаточные данные,
особенно конфигурацию
\key{две стороны и угол, не заключённый между ними}.

\end{itemize}

\subsection*{Главная логика доказательств}

В большинстве задач цель выглядит как равенство двух углов
или двух отрезков.

Полезный вопрос:

\[
\text{«В какие два треугольника входят эти элементы?»}
\]

Далее ищите один из трёх признаков равенства треугольников.

После доказательства равенства треугольников
используйте соответствующие элементы.

% ============================================================
% 1. УГЛЫ
% ============================================================

\section*{1. Углы при пересечении прямых}
\sectionline

\begin{minipage}[t]{0.50\linewidth}

\key{Вертикальные углы} равны:

\[
\angle 1=\angle 3,
\qquad
\angle 2=\angle 4.
\]

\key{Смежные углы} в сумме дают $180^\circ$:

\[
\angle 1+\angle 2=180^\circ.
\]

Если один угол равен $62^\circ$,
то вертикальный ему угол также равен $62^\circ$,
а каждый смежный равен

\[
180^\circ-62^\circ=118^\circ.
\]

\end{minipage}
\hfill
\begin{minipage}[t]{0.45\linewidth}

\centering

\begin{tikzpicture}[scale=0.88]

\tkzDefPoint(-2.3,-1.4){A}
\tkzDefPoint(2.3,1.4){B}
\tkzDefPoint(-2.3,1.4){C}
\tkzDefPoint(2.3,-1.4){D}

\tkzInterLL(A,B)(C,D)
\tkzGetPoint{O}

\tkzDrawSegments(A,B C,D)
\tkzDrawPoints(O)

\tkzMarkAngle[size=0.55,color=accent](B,O,C)
\tkzMarkAngle[size=0.55,color=accent](C,O,A)
\tkzMarkAngle[size=0.55,color=accent](A,O,D)
\tkzMarkAngle[size=0.55,color=accent](D,O,B)

\tkzLabelAngle[pos=0.8](B,O,C){$1$}
\tkzLabelAngle[pos=0.8](C,O,A){$2$}
\tkzLabelAngle[pos=0.8](A,O,D){$3$}
\tkzLabelAngle[pos=0.8](D,O,B){$4$}

\end{tikzpicture}

\end{minipage}

\subsection*{Ещё одна полезная конфигурация}

\begin{center}

\begin{tikzpicture}[scale=0.9]

\tkzDefPoint(-3,0){A}
\tkzDefPoint(3,0){B}
\tkzDefPoint(0,0){O}
\tkzDefPoint(1.8,2){C}

\tkzDrawSegments(A,B O,C)
\tkzDrawPoints(O)

\tkzMarkAngle[size=0.75,color=accent](A,O,C)
\tkzMarkAngle[size=0.55,color=accent](C,O,B)

\tkzLabelAngle[pos=1.02](A,O,C){$\alpha$}
\tkzLabelAngle[pos=0.82](C,O,B){$\beta$}

\end{tikzpicture}

\[
\alpha+\beta=180^\circ.
\]

\end{center}

\subsection*{Что проверить перед ответом}

\begin{itemize}

\item
\checkitem
Есть ли у углов общая вершина?

\item
\checkitem
Вертикальные ли это углы,
или они образуют развёрнутый угол?

\item
\checkitem
Правильно ли применено равенство вертикальных углов
или сумма смежных?

\end{itemize}

% ============================================================
% 2. ПАРАЛЛЕЛЬНЫЕ ПРЯМЫЕ
% ============================================================

\section*{2. Параллельные прямые и секущая}
\sectionline

Пусть

\[
a\parallel b,
\]

а прямая $c$ является секущей.

Тогда:

\begin{itemize}

\item
\key{накрест лежащие углы} равны;

\item
\key{соответственные углы} равны;

\item
сумма
\key{односторонних внутренних углов}
равна $180^\circ$.

\end{itemize}

\begin{center}

\begin{tikzpicture}[scale=0.92]

\tkzDefPoint(-3,1.3){A}
\tkzDefPoint(3,1.3){B}
\tkzDefPoint(-3,-1.3){C}
\tkzDefPoint(3,-1.3){D}

\tkzDefPoint(-1.1,-2.2){E}
\tkzDefPoint(1.1,2.2){F}

\tkzInterLL(A,B)(E,F)
\tkzGetPoint{P}

\tkzInterLL(C,D)(E,F)
\tkzGetPoint{Q}

\tkzDrawSegments(A,B C,D E,F)
\tkzDrawPoints(P,Q)

\tkzLabelPoint[above left](A){$a$}
\tkzLabelPoint[below left](C){$b$}
\tkzLabelPoint[right](F){$c$}

\tkzMarkAngle[size=0.45,color=accent](A,P,Q)
\tkzMarkAngle[size=0.45,color=accent](P,Q,D)

\end{tikzpicture}

\end{center}

На рисунке отмечена пара накрест лежащих углов.
При $a\parallel b$ они равны.

\subsection*{Три пары углов на одном рисунке}

\begin{center}

\begin{tikzpicture}[scale=0.9]

\tkzDefPoint(-3,1.2){A}
\tkzDefPoint(3,1.2){B}
\tkzDefPoint(-3,-1.2){C}
\tkzDefPoint(3,-1.2){D}

\tkzDefPoint(-0.8,-2.2){E}
\tkzDefPoint(1.0,2.2){F}

\tkzInterLL(A,B)(E,F)
\tkzGetPoint{P}

\tkzInterLL(C,D)(E,F)
\tkzGetPoint{Q}

\tkzDrawSegments(A,B C,D E,F)
\tkzDrawPoints(P,Q)

\tkzMarkAngle[size=0.40,color=accent](A,P,Q)
\tkzMarkAngle[size=0.40,color=accent](P,Q,D)

\tkzMarkAngle[size=0.58,color=darkaccent](F,P,B)
\tkzMarkAngle[size=0.58,color=darkaccent](P,Q,D)

\tkzLabelPoint[left](P){$P$}
\tkzLabelPoint[left](Q){$Q$}

\end{tikzpicture}

\end{center}

При решении сначала определите,
\key{какое взаимное положение занимают выбранные углы}.
Это снижает риск случайного применения неподходящего свойства.

\subsection*{Обратные признаки параллельности}

Если при пересечении двух прямых секущей выполняется одно из условий

\begin{itemize}

\item
накрест лежащие углы равны;

\item
соответственные углы равны;

\item
сумма односторонних внутренних углов
равна $180^\circ$,

\end{itemize}

то прямые параллельны.

\key{Типичная ошибка.}
Свойства углов при параллельных прямых применяйте,
когда параллельность дана или уже доказана.

% ============================================================
% 3. ТРЕУГОЛЬНИК
% ============================================================

\section*{3. Углы треугольника}
\sectionline

Для любого треугольника

\[
\angle A+\angle B+\angle C=180^\circ.
\]

\begin{center}

\begin{tikzpicture}[scale=0.88]

\tkzDefPoint(-2.5,0){A}
\tkzDefPoint(0.2,2.8){B}
\tkzDefPoint(2.5,0){C}

\tkzDrawPolygon(A,B,C)
\tkzDrawPoints(A,B,C)

\tkzLabelPoints[below](A,C)
\tkzLabelPoints[above](B)

\tkzMarkAngle[size=0.52,color=accent](C,A,B)
\tkzMarkAngle[size=0.48,color=accent](A,B,C)
\tkzMarkAngle[size=0.52,color=accent](B,C,A)

\end{tikzpicture}

\end{center}

Например, если

\[
\angle A=54^\circ,
\qquad
\angle B=71^\circ,
\]

то

\[
\angle C
=
180^\circ-54^\circ-71^\circ
=
55^\circ.
\]

% ============================================================
% 4. РАВНОБЕДРЕННЫЙ ТРЕУГОЛЬНИК
% ============================================================

\section*{4. Равнобедренный треугольник}
\sectionline

Если

\[
AB=BC,
\]

то треугольник $ABC$ равнобедренный
с основанием $AC$.

Следовательно,

\[
\angle A=\angle C.
\]

\begin{minipage}[t]{0.52\linewidth}

\subsection*{Пример}

Пусть

\[
AB=BC,
\qquad
\angle B=38^\circ.
\]

Тогда

\[
\angle A=\angle C.
\]

По сумме углов:

\[
\angle A+\angle C
=
180^\circ-38^\circ
=
142^\circ.
\]

Следовательно,

\[
\angle A=\angle C=71^\circ.
\]

\end{minipage}
\hfill
\begin{minipage}[t]{0.42\linewidth}

\centering

\begin{tikzpicture}[scale=0.82]

\tkzDefPoint(-2,0){A}
\tkzDefPoint(0,3){B}
\tkzDefPoint(2,0){C}

\tkzDrawPolygon(A,B,C)
\tkzDrawPoints(A,B,C)

\tkzLabelPoints[below](A,C)
\tkzLabelPoints[above](B)

\tkzMarkSegments[mark=|](A,B B,C)

\tkzMarkAngle[size=0.55,color=accent](A,B,C)
\tkzLabelAngle[pos=0.9](A,B,C){$38^\circ$}

\end{tikzpicture}

\end{minipage}

\subsection*{Медиана к основанию}

В равнобедренном треугольнике медиана,
проведённая из вершины к основанию,
одновременно является высотой и биссектрисой.

\begin{center}

\begin{tikzpicture}[scale=0.9]

\tkzDefPoint(-2.4,0){B}
\tkzDefPoint(0,3){A}
\tkzDefPoint(2.4,0){C}

\tkzDefMidPoint(B,C)
\tkzGetPoint{D}

\tkzDrawPolygon(A,B,C)
\tkzDrawSegment(A,D)
\tkzDrawPoints(A,B,C,D)

\tkzLabelPoints[below](B,C,D)
\tkzLabelPoints[above](A)

\tkzMarkSegments[mark=|](A,B A,C)
\tkzMarkSegments[mark=||](B,D D,C)

\tkzMarkAngle[size=0.48,color=accent](B,A,D)
\tkzMarkAngle[size=0.62,color=accent](D,A,C)

\tkzMarkRightAngle[size=0.3](A,D,C)

\end{tikzpicture}

\end{center}

Здесь одновременно выполняются:

\[
BD=DC,
\]

\[
\angle BAD=\angle DAC,
\]

\[
AD\perp BC.
\]

% ============================================================
% 5. ПРИЗНАКИ РАВЕНСТВА
% ============================================================

\section*{5. Три признака равенства треугольников}
\sectionline

\subsection*{I признак: СУС}

Если две стороны и угол между ними
одного треугольника соответственно равны
двум сторонам и углу между ними другого треугольника,
то треугольники равны.

\[
AB=DE,
\qquad
AC=DF,
\qquad
\angle BAC=\angle EDF.
\]

Тогда

\[
\triangle ABC\cong\triangle DEF.
\]

\begin{center}

\begin{tikzpicture}[scale=0.75]

\tkzDefPoint(-5,0){A}
\tkzDefPoint(-3.1,2.5){B}
\tkzDefPoint(-1.4,0){C}

\tkzDrawPolygon(A,B,C)
\tkzDrawPoints(A,B,C)

\tkzLabelPoints[below](A,C)
\tkzLabelPoints[above](B)

\tkzMarkSegments[mark=|](A,B)
\tkzMarkSegments[mark=||](A,C)
\tkzMarkAngle[size=0.55,color=accent](C,A,B)

\tkzDefPoint(1.4,0){D}
\tkzDefPoint(3.3,2.5){E}
\tkzDefPoint(5,0){F}

\tkzDrawPolygon(D,E,F)
\tkzDrawPoints(D,E,F)

\tkzLabelPoints[below](D,F)
\tkzLabelPoints[above](E)

\tkzMarkSegments[mark=|](D,E)
\tkzMarkSegments[mark=||](D,F)
\tkzMarkAngle[size=0.55,color=accent](F,D,E)

\end{tikzpicture}

\end{center}

\key{Контроль:}
угол обязательно расположен между двумя указанными сторонами.

\subsection*{II признак: УСУ}

Если сторона и два прилежащих к ней угла
одного треугольника соответственно равны
стороне и двум прилежащим к ней углам другого,
то треугольники равны.

\[
AB=DE,
\qquad
\angle A=\angle D,
\qquad
\angle B=\angle E.
\]

Тогда

\[
\triangle ABC\cong\triangle DEF.
\]

\begin{center}

\begin{tikzpicture}[scale=0.75]

\tkzDefPoint(-5,0){A}
\tkzDefPoint(-2.2,0){B}
\tkzDefPoint(-3.5,2.5){C}

\tkzDrawPolygon(A,B,C)
\tkzDrawPoints(A,B,C)
\tkzLabelPoints[below](A,B)
\tkzLabelPoints[above](C)

\tkzMarkSegments[mark=|](A,B)
\tkzMarkAngle[size=0.45,color=accent](B,A,C)
\tkzMarkAngle[size=0.55,color=accent](C,B,A)

\tkzDefPoint(1.8,0){D}
\tkzDefPoint(4.6,0){E}
\tkzDefPoint(3.3,2.5){F}

\tkzDrawPolygon(D,E,F)
\tkzDrawPoints(D,E,F)
\tkzLabelPoints[below](D,E)
\tkzLabelPoints[above](F)

\tkzMarkSegments[mark=|](D,E)
\tkzMarkAngle[size=0.45,color=accent](E,D,F)
\tkzMarkAngle[size=0.55,color=accent](F,E,D)

\end{tikzpicture}

\end{center}

\subsection*{III признак: ССС}

Если три стороны одного треугольника
соответственно равны трём сторонам другого,
то треугольники равны.

\[
AB=DE,
\qquad
BC=EF,
\qquad
AC=DF.
\]

Тогда

\[
\triangle ABC\cong\triangle DEF.
\]

\begin{center}

\begin{tikzpicture}[scale=0.75]

\tkzDefPoint(-5,0){A}
\tkzDefPoint(-3.4,2.5){B}
\tkzDefPoint(-1.2,0){C}

\tkzDrawPolygon(A,B,C)
\tkzDrawPoints(A,B,C)
\tkzLabelPoints[below](A,C)
\tkzLabelPoints[above](B)

\tkzMarkSegments[mark=|](A,B)
\tkzMarkSegments[mark=||](B,C)
\tkzMarkSegments[mark=|||](A,C)

\tkzDefPoint(1.2,0){D}
\tkzDefPoint(2.8,2.5){E}
\tkzDefPoint(5,0){F}

\tkzDrawPolygon(D,E,F)
\tkzDrawPoints(D,E,F)
\tkzLabelPoints[below](D,F)
\tkzLabelPoints[above](E)

\tkzMarkSegments[mark=|](D,E)
\tkzMarkSegments[mark=||](E,F)
\tkzMarkSegments[mark=|||](D,F)

\end{tikzpicture}

\end{center}

% ============================================================
% 6. СООТВЕТСТВИЕ ВЕРШИН
% ============================================================

\section*{6. Соответствие вершин}
\sectionline

Запись

\[
\triangle ABC\cong\triangle DEF
\]

означает:

\[
A\leftrightarrow D,
\qquad
B\leftrightarrow E,
\qquad
C\leftrightarrow F.
\]

Следовательно,

\[
AB=DE,
\qquad
BC=EF,
\qquad
AC=DF,
\]

а также

\[
\angle A=\angle D,
\qquad
\angle B=\angle E,
\qquad
\angle C=\angle F.
\]

\begin{center}

\begin{tikzpicture}[scale=0.72]

\tkzDefPoint(-5,0){A}
\tkzDefPoint(-3.4,2.5){B}
\tkzDefPoint(-1.2,0){C}

\tkzDrawPolygon(A,B,C)
\tkzDrawPoints(A,B,C)
\tkzLabelPoints[below](A,C)
\tkzLabelPoints[above](B)

\tkzDefPoint(1.2,0){D}
\tkzDefPoint(2.8,2.5){E}
\tkzDefPoint(5,0){F}

\tkzDrawPolygon(D,E,F)
\tkzDrawPoints(D,E,F)
\tkzLabelPoints[below](D,F)
\tkzLabelPoints[above](E)

\draw[
  -{Latex[length=2mm]},
  color=darkaccent,
  line width=0.7pt,
  dashed
]
(-4.7,-0.55)
to[bend right=13]
(1.45,-0.55);

\draw[
  -{Latex[length=2mm]},
  color=darkaccent,
  line width=0.7pt,
  dashed
]
(-3.4,3.0)
to[bend left=9]
(2.8,3.0);

\draw[
  -{Latex[length=2mm]},
  color=darkaccent,
  line width=0.7pt,
  dashed
]
(-1.0,-0.55)
to[bend right=13]
(4.8,-0.55);

\end{tikzpicture}

\end{center}

\subsection*{Недостаточные данные}

Если известно

\[
AB=DE,
\qquad
AC=DF,
\qquad
\angle B=\angle E,
\]

этого в общем случае недостаточно.

Угол $B$ не расположен между сторонами $AB$ и $AC$.

\begin{center}

\begin{tikzpicture}[scale=0.8]

\tkzDefPoint(-2.4,0){A}
\tkzDefPoint(0,2.6){B}
\tkzDefPoint(2.4,0){C}

\tkzDrawPolygon(A,B,C)
\tkzDrawPoints(A,B,C)

\tkzLabelPoints[below](A,C)
\tkzLabelPoints[above](B)

\tkzMarkSegments[mark=|](A,B)
\tkzMarkSegments[mark=||](A,C)

\tkzMarkAngle[size=0.55,color=accent](A,B,C)

\end{tikzpicture}

\end{center}

Конфигурация
«две стороны и произвольный угол»
в перечень трёх признаков равенства треугольников
не входит.

% ============================================================
% 7. ДОКАЗАТЕЛЬСТВА
% ============================================================

\section*{7. Типовые доказательства}
\sectionline

\subsection*{Пример 1. Диагональ параллелограмма}

В параллелограмме $ABCD$
проведена диагональ $AC$.

Докажем:

\[
\triangle ABC\cong\triangle CDA.
\]

У параллелограмма противоположные стороны равны:

\[
AB=CD,
\qquad
BC=AD.
\]

Сторона $AC$ общая:

\[
AC=CA.
\]

Получили три пары равных сторон, поэтому

\[
\triangle ABC\cong\triangle CDA
\]

по ССС.

\begin{center}

\begin{tikzpicture}[scale=0.84]

\tkzDefPoint(-2.8,0){A}
\tkzDefPoint(-1.2,2.0){B}
\tkzDefPoint(2.8,2.0){C}
\tkzDefPoint(1.2,0){D}

\tkzDrawPolygon(A,B,C,D)
\tkzDrawSegment(A,C)

\tkzDrawPoints(A,B,C,D)

\tkzLabelPoints[below](A,D)
\tkzLabelPoints[above](B,C)

\tkzMarkSegments[mark=|](A,B C,D)
\tkzMarkSegments[mark=||](B,C A,D)

\end{tikzpicture}

\end{center}

\subsection*{Пример 2. Четырёхугольник с двумя парами равных сторон}

Пусть

\[
AB=AD,
\qquad
CB=CD.
\]

Рассмотрим треугольники $ABC$ и $ADC$.

Имеем

\[
AB=AD,
\]

\[
BC=DC,
\]

\[
AC=AC.
\]

Следовательно,

\[
\triangle ABC\cong\triangle ADC
\]

по ССС.

Значит,

\[
\angle BAC=\angle CAD.
\]

Следовательно, $AC$ является биссектрисой угла $BAD$.

\begin{center}

\begin{tikzpicture}[scale=0.84]

\tkzDefPoint(-3,0){A}
\tkzDefPoint(0,1.8){B}
\tkzDefPoint(2.6,0){C}
\tkzDefPoint(0,-1.8){D}

\tkzDrawPolygon(A,B,C,D)
\tkzDrawSegment(A,C)

\tkzDrawPoints(A,B,C,D)

\tkzLabelPoints[left](A)
\tkzLabelPoints[above](B)
\tkzLabelPoints[right](C)
\tkzLabelPoints[below](D)

\tkzMarkSegments[mark=|](A,B A,D)
\tkzMarkSegments[mark=||](B,C D,C)

\tkzMarkAngle[size=0.45,color=accent](B,A,C)
\tkzMarkAngle[size=0.62,color=accent](C,A,D)

\end{tikzpicture}

\end{center}

\subsection*{Пример 3. Медиана к основанию}

Пусть

\[
AB=AC,
\]

точка $D$ лежит на $BC$ и

\[
BD=DC.
\]

Рассмотрим треугольники $ABD$ и $ACD$.

Имеем

\[
AB=AC,
\qquad
BD=DC,
\qquad
AD=AD.
\]

По ССС

\[
\triangle ABD\cong\triangle ACD.
\]

Следовательно,

\[
\angle BAD=\angle DAC.
\]

Значит, $AD$ является биссектрисой угла $A$.

\begin{center}

\begin{tikzpicture}[scale=0.84]

\tkzDefPoint(-2.3,0){B}
\tkzDefPoint(0,3){A}
\tkzDefPoint(2.3,0){C}

\tkzDefMidPoint(B,C)
\tkzGetPoint{D}

\tkzDrawPolygon(A,B,C)
\tkzDrawSegment(A,D)

\tkzDrawPoints(A,B,C,D)

\tkzLabelPoints[below](B,C,D)
\tkzLabelPoints[above](A)

\tkzMarkSegments[mark=|](A,B A,C)
\tkzMarkSegments[mark=||](B,D D,C)

\tkzMarkAngle[size=0.5,color=accent](B,A,D)
\tkzMarkAngle[size=0.65,color=accent](D,A,C)

\end{tikzpicture}

\end{center}

% ============================================================
% 8. ИСТОЧНИКИ РАВЕНСТВА
% ============================================================

\section*{8. Где искать недостающие равенства}
\sectionline

\begin{tabularx}
{\linewidth}
{@{}
>{\bfseries\raggedright\arraybackslash}
p{0.30\linewidth}
>
{\raggedright\arraybackslash}
X
@{}}

\toprule

Источник
&
Что можно получить
\\

\midrule

Общая сторона
&
Например,
$AC=AC$.
Это полноценная пара равных сторон.
\\

Вертикальные углы
&
Равенство углов при пересечении двух прямых.
\\

Параллельные прямые
&
Равные накрест лежащие
или соответственные углы.
\\

Равнобедренный треугольник
&
Равные углы при основании.
\\

Середина отрезка
&
Если $D$ --- середина $BC$, то

\[
BD=DC.
\]
\\

Биссектриса
&
Если $AD$ --- биссектриса $\angle A$, то

\[
\angle BAD=\angle DAC.
\]
\\

Прямой угол
&
Все прямые углы равны.
\\

\bottomrule

\end{tabularx}

\subsection*{Наглядная памятка}

\begin{center}

\begin{tikzpicture}[scale=0.72]

\tkzDefPoint(-6,0){A}
\tkzDefPoint(-4.4,2.3){B}
\tkzDefPoint(-2.8,0){C}
\tkzDefPoint(-4.4,-2.0){D}

\tkzDrawPolygon(A,B,C,D)
\tkzDrawSegment(A,C)

\tkzDrawPoints(A,C)
\tkzLabelPoints[left](A)
\tkzLabelPoints[right](C)

\tkzDefPoint(0,-1){E}
\tkzDefPoint(5,-1){F}
\tkzDefMidPoint(E,F)
\tkzGetPoint{M}

\tkzDrawSegment(E,F)
\tkzDrawPoints(E,F,M)
\tkzLabelPoints[below](E,F,M)
\tkzMarkSegments[mark=|](E,M M,F)

\end{tikzpicture}

\end{center}

Слева диагональ $AC$ является общей стороной
для двух треугольников.

Справа точка $M$ --- середина $EF$, поэтому

\[
EM=MF.
\]

% ============================================================
% 9. ПРЯМОУГОЛЬНЫЕ ТРЕУГОЛЬНИКИ
% ============================================================

\section*{9. Полезный приём с прямоугольными треугольниками}
\sectionline

В прямоугольном треугольнике
сумма двух острых углов равна $90^\circ$.

Если один острый угол равен $\alpha$, то второй:

\[
90^\circ-\alpha.
\]

\begin{center}

\begin{tikzpicture}[scale=0.86]

\tkzDefPoint(-2,0){A}
\tkzDefPoint(-2,2.8){B}
\tkzDefPoint(2.2,0){C}

\tkzDrawPolygon(A,B,C)
\tkzDrawPoints(A,B,C)

\tkzLabelPoints[left](A,B)
\tkzLabelPoints[below](C)

\tkzMarkRightAngle[size=0.35](B,A,C)
\tkzMarkAngle[size=0.55,color=accent](A,C,B)
\tkzLabelAngle[pos=0.9](A,C,B){$\alpha$}

\end{tikzpicture}

\end{center}

Если у двух прямоугольных треугольников
равны гипотенузы и по одному острому углу,
можно вычислить вторые острые углы
и затем применить УСУ.

% ============================================================
% АЛГОРИТМ
% ============================================================

\clearpage

\thispagestyle{plain}

\begin{center}

{\Large
\textcolor{darkaccent}{
\textbf{
Алгоритм: доказательство равенства отрезков или углов
}
}
}

\end{center}

\vspace{7mm}

\begin{center}

\begin{tikzpicture}[
  node distance=8mm and 18mm
]

\node[flowstart](start)
{
Начало
};

\node[
  flowact,
  below=of start
]
(goal)
{
Определите,
какие два отрезка или угла
требуется сравнить
};

\node[
  flowact,
  below=of goal
]
(tri)
{
Найдите два треугольника,
содержащие сравниваемые элементы
};

\node[
  flowact,
  below=of tri
]
(facts)
{
Соберите равные элементы:
данные, общая сторона,
вертикальные углы,
параллельность,
равнобедренность
};

\node[
  flowcond,
  below=11mm of facts
]
(criterion)
{
Есть СУС,
УСУ или ССС?
};

\node[
  flowact,
  below=11mm of criterion
]
(equal)
{
Запишите равенство треугольников
в правильном порядке вершин
};

\node[
  flowact,
  below=of equal
]
(concl)
{
Используйте равенство
соответствующих сторон или углов
};

\node[
  flowstart,
  below=of concl
]
(finish)
{
Готово
};

\node[
  flowact,
  right=25mm of criterion,
  text width=46mm
]
(search)
{
Получите ещё одно равенство
или выберите другую пару треугольников
};

\draw[flowarr](start)--(goal);
\draw[flowarr](goal)--(tri);
\draw[flowarr](tri)--(facts);
\draw[flowarr](facts)--(criterion);

\draw[flowarr]
(criterion)
--
node[flowlabel,left]{да}
(equal);

\draw[flowarr](equal)--(concl);
\draw[flowarr](concl)--(finish);

\draw[flowarr]
(criterion.east)
--
node[flowlabel,above]{нет}
(search.west);

\draw[flowarr]
(search.north)
|-
([xshift=8mm]facts.east)
--
(facts.east);

\end{tikzpicture}

\end{center}

\clearpage

% ============================================================
% 10. ТИПИЧНЫЕ ОШИБКИ
% ============================================================

\section*{10. Типичные ошибки МЦКО}
\sectionline

\begin{enumerate}[label=\arabic*)]

\item
\key{Вертикальные углы.}

Для них записывают сумму $180^\circ$,
хотя требуется использовать равенство.

\item
\key{Смежные углы.}

Для них записывают равенство,
хотя требуется сумма $180^\circ$.

\item
\key{Параллельные прямые.}

Используют свойства углов,
хотя параллельность ещё не установлена.

\item
\key{СУС.}

Выбирают угол,
который не заключён между двумя указанными сторонами.

\item
\key{Порядок вершин.}

Записывают равенство треугольников
в несогласованном порядке.

\item
\key{Равнобедренный треугольник.}

Путают основание и боковые стороны.

\item
\key{Доказательство.}

Требуемое равенство используют в качестве исходного факта.

\item
\key{Чертёж.}

Делают вывод только по внешнему виду рисунка.
Каждый факт должен иметь основание.

\end{enumerate}

\subsection*{Мини-проверка перед сдачей}

\begin{itemize}

\item
\checkitem
Я указал источник каждого равенства.

\item
\checkitem
Я назвал конкретный признак
равенства треугольников.

\item
\checkitem
Я проверил соответствие вершин.

\item
\checkitem
Финальный вывод отвечает
именно на вопрос задачи.

\end{itemize}

% ============================================================
% ТРЕНИРОВОЧНЫЙ БЛОК
% ============================================================

\clearpage

\section*{Тренировочный блок · Путь А}
\sectionline

Задания расположены по нарастающей сложности.

Решения в тренировочном блоке не приводятся.

\subsection*{Средний уровень}

\begin{enumerate}[label=\textbf{\arabic*.}]

\item
Две прямые пересекаются.
Один из образовавшихся углов равен

\[
47^\circ.
\]

Найдите остальные три угла.

\begin{center}

\begin{tikzpicture}[scale=0.62]

\tkzDefPoint(-2,-1.1){A}
\tkzDefPoint(2,1.1){B}
\tkzDefPoint(-2,1.1){C}
\tkzDefPoint(2,-1.1){D}

\tkzInterLL(A,B)(C,D)
\tkzGetPoint{O}

\tkzDrawSegments(A,B C,D)
\tkzMarkAngle[size=0.5,color=accent](B,O,C)
\tkzLabelAngle[pos=0.78](B,O,C){$47^\circ$}

\end{tikzpicture}

\end{center}

\item
Прямые

\[
a\parallel b
\]

пересечены секущей $c$.

Один из внутренних односторонних углов
равен

\[
116^\circ.
\]

Найдите второй внутренний односторонний угол
и равный ему соответственный острый угол.

\item
В равнобедренном треугольнике $ABC$
известно:

\[
AB=BC,
\qquad
\angle B=46^\circ.
\]

Найдите углы $A$ и $C$.

\begin{center}

\begin{tikzpicture}[scale=0.62]

\tkzDefPoint(-2,0){A}
\tkzDefPoint(0,2.6){B}
\tkzDefPoint(2,0){C}

\tkzDrawPolygon(A,B,C)
\tkzMarkSegments[mark=|](A,B B,C)

\tkzLabelPoints[below](A,C)
\tkzLabelPoints[above](B)

\tkzMarkAngle[size=0.52,color=accent](A,B,C)
\tkzLabelAngle[pos=0.82](A,B,C){$46^\circ$}

\end{tikzpicture}

\end{center}

\end{enumerate}

\subsection*{Выше среднего}

\begin{enumerate}[
  label=\textbf{\arabic*.},
  start=4
]

\item
В треугольниках $ABC$ и $DEF$ известно:

\[
AB=DE,
\qquad
AC=DF,
\qquad
\angle BAC=\angle EDF.
\]

Докажите равенство треугольников,
назовите признак
и укажите соответствие вершин.

\item
В треугольниках $ABC$ и $DEF$ известно:

\[
AB=DE,
\qquad
AC=DF,
\qquad
\angle B=\angle E.
\]

Можно ли только по этим данным
гарантировать равенство треугольников?

Обоснуйте ответ.

\item
Дан параллелограмм $ABCD$.

Проведена диагональ $AC$.

Докажите равенство треугольников

\[
ABC
\]

и

\[
CDA.
\]

Укажите признак.

\begin{center}

\begin{tikzpicture}[scale=0.62]

\tkzDefPoint(-2.5,0){A}
\tkzDefPoint(-1.2,1.7){B}
\tkzDefPoint(2.5,1.7){C}
\tkzDefPoint(1.2,0){D}

\tkzDrawPolygon(A,B,C,D)
\tkzDrawSegment(A,C)

\tkzLabelPoints[below](A,D)
\tkzLabelPoints[above](B,C)

\end{tikzpicture}

\end{center}

\item
В четырёхугольнике $ABCD$ известно:

\[
AB=AD,
\qquad
CB=CD.
\]

Проведена диагональ $AC$.

Докажите:

\[
\angle BAC=\angle CAD.
\]

\item
В равнобедренном треугольнике $ABC$ дано

\[
AB=AC.
\]

Точка $D$ --- середина основания $BC$.

Докажите,
что $AD$ является биссектрисой угла $A$.

\item
В треугольнике $ABC$
точки $M$ и $N$ лежат соответственно
на сторонах $AB$ и $AC$.

Известно:

\[
AB=AC,
\qquad
AM=AN.
\]

Докажите:

\[
BN=CM.
\]

\begin{center}

\begin{tikzpicture}[scale=0.7]

\tkzDefPoint(-2.6,0){B}
\tkzDefPoint(0,3){A}
\tkzDefPoint(2.6,0){C}

\tkzDefPoint(-1.05,1.8){M}
\tkzDefPoint(1.05,1.8){N}

\tkzDrawPolygon(A,B,C)
\tkzDrawSegment(B,N)
\tkzDrawSegment(C,M)

\tkzDrawPoints(A,B,C,M,N)

\tkzLabelPoints[below](B,C)
\tkzLabelPoints[above](A)
\tkzLabelPoint[left](M){$M$}
\tkzLabelPoint[right](N){$N$}

\end{tikzpicture}

\end{center}

\end{enumerate}

\subsection*{Сложный уровень}

\begin{enumerate}[
  label=\textbf{\arabic*.},
  start=10
]

\item
В равнобедренном треугольнике $ABC$
известно

\[
AB=AC.
\]

Точка $D$ --- середина $BC$,
а точка $E$ лежит на отрезке $AD$.

Докажите:

\[
BE=CE.
\]

Постройте доказательство в два этапа:

сначала получите равенство углов при вершине $A$,
затем выберите новую пару треугольников.

\begin{center}

\begin{tikzpicture}[scale=0.72]

\tkzDefPoint(-2.5,0){B}
\tkzDefPoint(0,3.2){A}
\tkzDefPoint(2.5,0){C}

\tkzDefMidPoint(B,C)
\tkzGetPoint{D}

\tkzDefPoint(0,1.4){E}

\tkzDrawPolygon(A,B,C)
\tkzDrawSegment(A,D)
\tkzDrawSegments(B,E C,E)

\tkzDrawPoints(A,B,C,D,E)

\tkzLabelPoints[below](B,C,D)
\tkzLabelPoints[above](A)
\tkzLabelPoint[right](E){$E$}

\tkzMarkSegments[mark=|](A,B A,C)
\tkzMarkSegments[mark=||](B,D D,C)

\end{tikzpicture}

\end{center}

\end{enumerate}

% ============================================================
% ОТВЕТЫ
% ============================================================

\clearpage

\section*{Ответы к тренировочному блоку}
\sectionline

\footnotesize

\begin{longtable}
{@{}
p{0.07\linewidth}
p{0.86\linewidth}
@{}}

\toprule

\textbf{№}
&
\textbf{Ответ / ключевой факт}
\\

\midrule
\endfirsthead

\multicolumn{2}{l}{
\small\textcolor{softgray}{Продолжение таблицы ответов}
}
\\[2mm]

\toprule

\textbf{№}
&
\textbf{Ответ / ключевой факт}
\\

\midrule
\endhead

\midrule
\multicolumn{2}{r}{
\small\textcolor{softgray}{Продолжение на следующей странице}
}
\\
\endfoot

\bottomrule
\endlastfoot

1
&
Углы:

\[
47^\circ,\quad
133^\circ,\quad
47^\circ,\quad
133^\circ.
\]

Используются свойства вертикальных
и смежных углов.
\\[2mm]

2
&
Второй односторонний угол:

\[
180^\circ-116^\circ=64^\circ.
\]

Соответственный ему острый угол:

\[
64^\circ.
\]
\\[2mm]

3
&
Поскольку

\[
AB=BC,
\]

то

\[
\angle A=\angle C.
\]

Поэтому

\[
\angle A
=
\angle C
=
\frac{180^\circ-46^\circ}{2}
=
67^\circ.
\]
\\[2mm]

4
&
Имеем

\[
AB=DE,
\qquad
AC=DF,
\qquad
\angle BAC=\angle EDF.
\]

Следовательно,

\[
\triangle ABC\cong\triangle DEF
\]

по СУС.

Соответствие:

\[
A\leftrightarrow D,
\qquad
B\leftrightarrow E,
\qquad
C\leftrightarrow F.
\]
\\[2mm]

5
&
Гарантировать равенство нельзя.

Даны две стороны и угол,
который не расположен между этими сторонами.

Такой набор данных не совпадает
ни с СУС, ни с УСУ, ни с ССС.
\\[2mm]

6
&
В параллелограмме

\[
AB=CD,
\qquad
BC=AD.
\]

Кроме того,

\[
AC=CA.
\]

Следовательно,

\[
\triangle ABC\cong\triangle CDA
\]

по ССС.
\\[2mm]

7
&
Имеем

\[
AB=AD,
\qquad
BC=DC,
\qquad
AC=AC.
\]

Следовательно,

\[
\triangle ABC\cong\triangle ADC
\]

по ССС.

Отсюда

\[
\angle BAC=\angle CAD.
\]
\\[2mm]

8
&
Имеем

\[
AB=AC,
\qquad
BD=DC,
\qquad
AD=AD.
\]

Следовательно,

\[
\triangle ABD\cong\triangle ACD
\]

по ССС.

Поэтому

\[
\angle BAD=\angle DAC,
\]

то есть $AD$ --- биссектриса угла $A$.
\\[2mm]

9
&
Рассмотрим

\[
\triangle ABN
\]

и

\[
\triangle ACM.
\]

Имеем

\[
AB=AC,
\qquad
AN=AM.
\]

Поскольку $N\in AC$ и $M\in AB$,

\[
\angle BAN=\angle CAM=\angle BAC.
\]

Следовательно,

\[
\triangle ABN\cong\triangle ACM
\]

по СУС.

Отсюда

\[
BN=CM.
\]
\\[2mm]

10
&
Сначала рассмотрим треугольники

\[
ABD
\]

и

\[
ACD.
\]

Имеем

\[
AB=AC,
\qquad
BD=DC,
\qquad
AD=AD.
\]

По ССС:

\[
\triangle ABD\cong\triangle ACD.
\]

Следовательно,

\[
\angle BAD=\angle DAC.
\]

Поскольку $E\in AD$,

\[
\angle BAE=\angle EAC.
\]

Теперь рассмотрим

\[
\triangle ABE
\]

и

\[
\triangle ACE.
\]

Для них

\[
AB=AC,
\qquad
AE=AE,
\qquad
\angle BAE=\angle EAC.
\]

По СУС:

\[
\triangle ABE\cong\triangle ACE.
\]

Следовательно,

\[
BE=CE.
\]
\\

\end{longtable}

\normalsize

\end{document}
